% !PRINTED 498
\weberpage{498}
%<<CONT
und man kann also, wenn ausser den vier Ziffern 0, 1, 2, 3 noch eine fünfte, 4, vorhanden ist, alle dreigliedrigen Cyklen aus Transpositionspaaren von der Form $(0, 1)(3, 4)$ zusammensetzen.

Hieran schliessen wir die Beweise von zwei weiteren wichtigen Sätzen, die sich auf transitive Permutationsgruppen beziehen:

\gesperrt{9. Wenn eine transitive Permutationsgruppe von $m$ Ziffern eine einzelne Transposition enthält, so ist sie entweder die symmetrische Gruppe oder sie ist imprimitiv.}

Die vorgelegte Gruppe sei $P$, die $m$ Ziffern, die durch sie permutirt werden, $0, 1, 2 \dots m - 1$, und es komme darin die Transposition $(0, 1)$ vor. Sie möge überhaupt die Transpositionen
\[
(0, 1), (0, 2), (0, 3), \dots (0, \mu - 1),
\]
aber keine andere Transposition mit der Ziffer 0 enthalten. Ist nun $\mu = m$, so ist $P$ nach 2. die symmetrische Gruppe. Ist aber $\mu < m$, so enthält $P$ nach demselben Satze alle Permutationen der Ziffern
\[
M = 0, 1, 2 \dots \mu - 1.
\]
Daraus folgt, dass $P$ keine Transposition von einer der Ziffern $0, 1, 2, 3 \dots \mu - 1$ mit einer anderen Ziffer, etwa mit $\mu$, enthalten kann. Denn wenn $(1, \mu)$ in $P$ vorkommt, so kommt auch $(0, \mu) = (0, 1)(1, \mu)(0, 1)$ darin vor, gegen die Voraussetzung.

Wenn nun $P$ transitiv ist, so giebt es darin eine Permutation $\pi$, durch die die Ziffer 0 in eine nicht in dem System $M$ enthaltene Ziffer, etwa in $\mu$ übergeht, und dann kann durch $\pi$ keine der Ziffern von $M$ in eine Ziffer von $M$ übergehen; denn wenn z.\,B. 1 in $r$ übergeht, so ist
\[
\pi^{-1}(0, 1)\pi = (\mu, r)
\]
und $(\mu, r)$ kann, wenn $r$ zu $M$ gehört, wie wir gesehen haben, nicht in $P$ vorkommen. Durch jede Permutation von $P$ werden also die Ziffern des Systems $M$ entweder nur unter sich vertauscht, oder sie werden in die Ziffern eines ganz davon verschiedenen Systems $M'$ übergeführt, und in $P$ kommen auch alle Permutationen der Ziffern $M'$ vor.

Wenn mit $M$ und $M'$ noch nicht alle Ziffern erschöpft sind, so giebt es eine Permutation in $P$, durch die $M$ in ein drittes System $M''$ übergeführt wird, was wieder mit $M$ und mit $M'$ keine Ziffer gemein haben kann.
